\subsection{Pipeline sin fuga de información}

Todo lo que se aprende de los datos se ajusta solo con el entrenamiento de cada ventana y se aplica sin cambios a la validación o a la prueba:
\begin{itemize}
\item la transformación $\log(1+x)$ y la estandarización de la regresión logística, dentro de un \texttt{Pipeline};
\item la función de distribución empírica del modelo difuso;
\item el umbral de decisión.
\end{itemize}
No se balancean las clases: el AUC-PR se mide con la prevalencia real, y los conjuntos balanceados de forma artificial inflan el desempeño de manera engañosa \cite{Nguyen2025}. Las variables solo usan información anterior a la integración del commit. La experiencia del autor se cuenta con los commits previos en orden temporal.

\subsection{Partición temporal}

Nunca se mezclan commits al azar (Fig.~\ref{fig:particion}). Para cada año de validación $V\in\{2023,2024,2025\}$ se entrena con $[\text{2019-01-01},\ \text{inicio}(V)-90\text{ d})$ y se valida con el año $V$ completo. El hueco de 90 días replica lo que ocurre en la práctica: al entrenar, las etiquetas de los últimos 90 días todavía no existen. El modelo final se entrena con $[\text{2019-01-01},\ \text{2025-10-03})$ y se evalúa una sola vez con los \cifNPrueba{} commits de enero a julio de 2026 (prevalencia \cifPrevPrueba{}).

La partición de prueba quedó guardada con una huella SHA-256 de sus hashes y etiquetas. La evaluación final verifica de forma estricta que la huella no cambie. Se declara una salvedad: durante una fase preliminar de factibilidad, modelos simples se evaluaron una vez sobre 2026 para confirmar que había señal. Ninguna decisión posterior usó esas cifras.

\begin{figure}[!t]
\centering
% Partición temporal: ventana creciente con hueco de 90 días. Escala: 1 año = 0,85 cm desde 2019.
\begin{tikzpicture}[x=0.85cm,y=0.42cm,font=\scriptsize]
  \def\h{0.25}  % hueco de 90 días ~ 0,25 años
  \foreach \y/\lab/\ini in {3/V$_{2023}$/4, 2/V$_{2024}$/5, 1/V$_{2025}$/6}{
    \fill[udazul!70] (0,\y-0.32) rectangle (\ini-\h,\y+0.32);
    \fill[udplata!35] (\ini-\h,\y-0.32) rectangle (\ini,\y+0.32);
    \fill[uddorado!85] (\ini,\y-0.32) rectangle (\ini+1,\y+0.32);
    \node[left,udsable] at (0,\y) {\lab};
  }
  \fill[udazul!70] (0,-0.32) rectangle (7-\h,0.32);
  \fill[udplata!35] (7-\h,-0.32) rectangle (7,0.32);
  \fill[udrojo!80] (7,-0.32) rectangle (7+7/12,0.32);
  \node[left,udsable] at (0,0) {Prueba};
  \draw[udsable] (0,-0.75) -- (8,-0.75);
  \foreach \x/\t in {0/2019,1/2020,2/2021,3/2022,4/2023,5/2024,6/2025,7/2026}
    \draw[udsable] (\x,-0.75) -- ++(0,-0.12) node[below] {\t};
  \node[anchor=west,fill=udazul!70,minimum width=3mm,minimum height=2mm,inner sep=0] (a) at (0,4.2) {};
  \node[anchor=west] at (a.east) {entrenamiento};
  \node[anchor=west,fill=udplata!35,minimum width=3mm,minimum height=2mm,inner sep=0] (b) at (2.2,4.2) {};
  \node[anchor=west] at (b.east) {hueco 90 d};
  \node[anchor=west,fill=uddorado!85,minimum width=3mm,minimum height=2mm,inner sep=0] (c) at (4.1,4.2) {};
  \node[anchor=west] at (c.east) {validación};
  \node[anchor=west,fill=udrojo!80,minimum width=3mm,minimum height=2mm,inner sep=0] (d) at (5.9,4.2) {};
  \node[anchor=west] at (d.east) {prueba};
\end{tikzpicture}

\caption{Partición temporal con ventana creciente. Cada ventana de validación (2023, 2024 y 2025) se entrena con todo lo anterior menos un hueco de 90 días; la prueba (enero--julio de 2026) se evalúa una sola vez.}
\label{fig:particion}
\end{figure}

\subsection{Métricas y pruebas estadísticas}

\textbf{Métrica principal: AUC-PR} (precisión promedio de la clase riesgoso). Con una prevalencia cercana al 4\,\%, el AUC-ROC es engañosamente alto y el AUC-PR muestra la utilidad real de la lista de commits marcados \cite{Saito2015}. El AUC-PR de un clasificador sin información es igual a la prevalencia; por eso también se reporta el cociente AUC-PR/prevalencia.

\textbf{Métricas secundarias:} MCC, AUC-ROC, sensibilidad y especificidad. Las métricas que dependen de un umbral usan el que maximiza el MCC en la ventana de validación 2025, aplicado tal cual en la prueba.

\textbf{Costo:} tiempo de entrenamiento, tiempo de inferencia y número de parámetros.

\textbf{Incertidumbre.} Se reportan intervalos de confianza del 95\,\% por \emph{bootstrap} (1000 remuestreos de la prueba) y media $\pm$ desviación en 5 semillas (11, 23, 37, 41 y 53) para los modelos estocásticos.

\textbf{Pruebas.}
\begin{itemize}
\item McNemar con corrección de continuidad, entre el difuso-evolutivo y cada referencia, sobre las mismas predicciones de prueba \cite{Dietterich1998}.
\item Friedman con \emph{post hoc} de Nemenyi sobre los rangos de AUC-PR de los cinco modelos en las tres ventanas de validación más la prueba \cite{Demsar2006}.
\end{itemize}
En ambos casos $\alpha=0{,}05$. Con solo 5 semillas no se usa Wilcoxon. Una diferencia sin prueba se reporta como ``diferencia observada''.

\subsection{Resultados de validación de las referencias}

Las referencias se ajustaron con el mismo pipeline y la misma partición que el modelo propuesto. Se probó una pequeña rejilla (\cifNConfigsRef{} configuraciones) y se eligió la de mayor AUC-PR medio en las tres ventanas.

\textbf{Mejor referencia clásica:} el \emph{gradient boosting}, con AUC-PR medio de \cifValGB{} y alrededor de \cifParamGB{} parámetros (nodos de los árboles). Es la referencia $R$ del criterio de éxito.

\textbf{Regresión logística:} da \cifValLogistica{} con \cifParamLR{} parámetros, y la regularización no cambia nada porque hay cientos de miles de ejemplos.

\textbf{Clase mayoritaria:} queda en \cifValMayoritaria{}, la prevalencia. Ambas referencias clásicas la multiplican por cerca de cuatro.


\subsection{Espacio de hiperparámetros del modelo propuesto}

El espacio se declara antes de buscar (Tabla~\ref{tab:espacio}). La búsqueda es aleatoria, con 14 configuraciones de las 972 posibles. Cada configuración se evalúa en las tres ventanas con la semilla 11 y se elige por el AUC-PR medio. La sensibilidad se analiza sobre la rejilla de máximo de reglas $\times$ probabilidad de mutación, con el resto fijo en la mejor configuración. Después, la configuración elegida se entrena con las 5 semillas.

\begin{table}[!t]
\caption{Espacio de hiperparámetros del difuso-evolutivo}
\label{tab:espacio}
\centering\footnotesize
\begin{tabularx}{\columnwidth}{@{}l Y@{}}
\toprule
\cabecera \textbf{Hiperparámetro} & \textbf{Valores}\\
\midrule
\filas{../comun/tabla_espacio.tex}
\bottomrule
\end{tabularx}
\end{table}
